% Working document for the DKRZ Levante application. Build: tectonic dkrz_antrag.tex
\documentclass[10pt,a4paper]{article}
\usepackage[margin=2.2cm]{geometry}
\usepackage{amsmath,booktabs,microtype}
\usepackage{enumitem}
\setlist{itemsep=2pt,topsep=3pt}
\usepackage[hidelinks]{hyperref}
\newcommand{\dx}{\Delta x}

\title{\vspace{-1.5em}DKRZ Levante application --- 20\,m WRF-LES of the Radfeld transect\\
\large Working document: core bullets, arithmetic, process constraints, proposal outline}
\author{Elias Wahl (Innsbruck, FWF) --- with the Tübingen UAS group (DFG)}
\date{2026-09-08}

\begin{document}
\maketitle

\section*{The five core bullets (agreed version: 50\,000 node-hours)}

\paragraph{1 --- Objective.} 20\,m WRF-LES of the Radfeld transect (Inn Valley,
TEAMx) to derive correction functions for UAS-measured turbulence --- the
objective of the DFG-funded partner project (virtual flights along the real
ascent trajectories, identical processing chain). Complementary objective of
the FWF-funded applicant: the same runs as turbulence reference for evaluating
a 3D PBL closure in the 500\,m grey zone.

\paragraph{2 --- Funding context.} The UAS data, processing chain and joint
analysis are contributed by the DFG-funded group at Tübingen (Platis); the
applicant's own position is FWF-funded. Results feed the DFG project's work
packages directly; research stay in Tübingen spring 2027.

\paragraph{3 --- Feasibility.} Production only: a validated ICON-driven 500\,m
WRF setup of the valley is in continuous research use with measured
throughput; a 100\,m LES of the same valley exists in the group (Freddi 2026)
and supplies the intermediate nest; no methodological development on Levante.

\paragraph{4 --- Scope and request.} Six full diurnal cycles from the two
TEAMx intensive periods --- nights included, because the valley cold pool and
stable-layer turbulence are where the model deficits concentrate and where
20\,m resolution is decisive; the 20\,m nest covers the full cross-valley
transect with along-valley fetch ($\sim$130$\times$10$^6$ points,
$\Delta t \approx 0.1$\,s, $\approx$160 node-h per simulated hour, anchored in
the measured 500\,m throughput). With a three-member sensitivity suite
(subgrid closure, nesting ratio) and a 50\,\% restart/I-O margin:
\textbf{50\,000 node-hours, $\approx$60\,TB work / 15\,TB archive}.

\paragraph{5 --- Data value.} One reference dataset serving three uses --- UAS
correction functions (DFG partner project), grey-zone evaluation of a 3D PBL
closure (FWF-funded applicant project), and the first LES-resolved nocturnal
cold-pool budget of this valley --- documented and shared with the TEAMx
community.

\section*{The arithmetic (for the resource section)}

\begin{center}
\begin{tabular}{@{}llr@{}}
\toprule
component & assumption & cost \\
\midrule
500\,m parent & 24$\times$10$^6$ pts, $\Delta t$ = 2\,s --- \textbf{measured} & $\approx$190 core-h / sim-h \\
100\,m nest & $\sim$40$\times$30\,km, 12$\times$10$^6$ pts, $\Delta t$ = 0.6\,s (Freddi-like) & $\approx$320 core-h / sim-h \\
20\,m nest & transect + fetch $\approx$ 25$\times$15\,km, 130$\times$10$^6$ pts, $\Delta t$ = 0.1\,s & $\approx$20\,000 core-h / sim-h \\
chain total & & $\approx$160 node-h / sim-h \\
\midrule
production & 6 cases $\times$ 26 sim-h (24\,h + spin-up) & $\approx$25\,000 node-h \\
sensitivity & 3 members $\times$ 1 case & $\approx$12\,000 node-h \\
margin & restarts, I/O, failed segments ($\times$1.33) & $\rightarrow$ $\approx$\textbf{50\,000 node-h} \\
pre/post & virtual-flight sampling, spectra, statistics (SLURM) & + 3\,000 node-h \\
\bottomrule
\end{tabular}
\end{center}

Fallback floor: 42\,k (two sensitivity members). Shelved 100\,k extensions:
10 cases, a 10\,m stable-night pilot ($\sim$5\,k) --- reusable as a follow-up
request. Scalability, measured on the 500\,m domain (128-core nodes):
1.5\,s/step on 2 nodes, $\sim$0.8 on 4, 0.65 on 5 --- near-linear to 5 nodes.

\textbf{To confirm before writing:} Freddi's 100\,m domain extent and time
step; the 20\,m nest extent (25$\times$15\,km assumed); whether one case needs
the multi-layer-stratification rehearsal.

\section*{DKRZ process constraints (allocation rules, 2026-09-08)}

\begin{itemize}
\item \textbf{Eligibility / PI:} the PI should be at a German research
  institution; international groups need a significant German contribution.
  $\Rightarrow$ the proposal is \textbf{led by Tübingen (Platis)}; Innsbruck
  joins as international partner supplying and running the model system.
\item \textbf{Windows:} Sep 1 -- Oct 31 (start Jan 1) and Mar 1 -- Apr 30
  (start Jul 1); allocations run 12 months. \textbf{Current window open,
  deadline 2026-10-31, start 2027-01-01.}
\item \textbf{Compute expires quarterly} --- consumption must be roughly even.
  A Jan 1 start forces $\sim$12.5\,k node-h per quarter from January; LES
  production is planned H2 2027, so either front-load the sensitivity suite
  into Q1/Q2 or target the spring window (start 2027-07-01).
  \emph{Decide before writing.}
\item \textbf{Pre-/post-processing} runs through SLURM and needs its own
  requested compute ($\approx$3\,k node-h, in the table).
\item \textbf{Document limits:} first proposal $\le$ 8 pages,
  $\le$ 4\,000 characters/page; required sections incl.\ scalability
  (covered by the measured scaling curve). Renewals: $\le$ 2-page report.
\item \textbf{No data-storage usage plan triggered} (thresholds:
  $\ge$ 10$^6$ node-h, $\ge$ 1\,024\,TiB disk or tape --- ours: 53\,k / 60\,TB
  / 15\,TB).
\item \textbf{Storage policy:} disk per allocation period (re-request on
  renewal; move data to archive within the period); unused archive grants
  expire --- request archive close to what will be written.
\item \textbf{GPU:} not requested (CPU-only WRF; only hardware-accelerated
  remote visualization would justify it).
\end{itemize}

\section*{Proposal outline ($\le$ 8 pages, DKRZ's section order)}

\begin{enumerate}
\item \textbf{Project overview} ($\sim$3\,500 chars): the TEAMx campaign and
  the UAS turbulence gap; the 20\,m LES as the closing reference; partners,
  funding, deliverables, data sharing.
\item \textbf{Planned work, scientific view} ($\sim$5\,000): six diurnal
  cycles (nights included); virtual UAS ascents $\rightarrow$ transfer
  functions vs.\ $\dx/z_i$; secondary uses (grey-zone reference, cold-pool
  budget). Only what the compute buys --- no closure-development story.
\item \textbf{Mathematical/computational aspects} ($\sim$4\,000): WRF 4.8;
  online nesting 500 $\rightarrow$ 100 $\rightarrow$ 20\,m; Deardorff TKE
  closure in the LES domains; RK3 split-explicit dynamics; ICON 500\,m
  forcing (native 65-level conversion); output design for virtual flights.
\item \textbf{Numerical methods} ($\sim$3\,000): time-step ratios
  (2 / 0.6 / 0.1\,s), nesting-ratio and damping choices with complex-terrain
  precedents (METEX21, Perdig\~ao, Riviera); absorbing layer; spin-up;
  the 3-member sensitivity suite.
\item \textbf{Suitability for HLRE-4} ($\sim$2\,500): pure-MPI WRF on Milan
  nodes; production, not development; group precedents (the running 500\,m
  system, the existing 100\,m LES); I/O sized to Lustre.
\item \textbf{Scalability} ($\sim$2\,500 + table): the measured 2/4/5-node
  curve; planned node counts per domain (10--20 nodes per LES segment).
\item \textbf{Computing time and storage} ($\sim$4\,500 + the arithmetic
  table): 50\,k + 3\,k node-h, 60\,TB work / 15\,TB archive; per-quarter
  consumption profile matched to the chosen start date.
\item \textbf{Added value} ($\sim$2\,500): one dataset, three uses; TEAMx
  reuse; no comparable LES at the transect.
\end{enumerate}

\section*{Open decisions}
\begin{enumerate}
\item Window: 2026-10-31 deadline (start Jan 1) vs.\ spring (start Jul 1) ---
  Elias + Platis, against the quarterly-expiry profile.
\item PI arrangement and writing split with Tübingen.
\item The three numbers to confirm (Freddi extent/time step, 20\,m nest
  extent, rehearsal case).
\end{enumerate}

\end{document}
